  \noindent\hfill
 \begin{minipage}[t]{0.6\textwidth}
 	%RESUMEN DE LA TESIS PARA OPTAR AL \\
 	%GRADO DE DOCTOR[A] EN INGENIERÍA \\
 	%ELÉCTRICA, \\
 	RESUMEN DE LA TESIS PARA OPTAR AL \\
 	GRADO DE MAGÍSTER EN CIENCIAS DE LA \\ INGENIERÍA, MENCIÓN ELÉCTRICA, \\
 	MEMORIA PARA OPTAR AL TÍTULO DE \\ INGENIERO CIVIL ELÉCTRICO \\
 	{POR: IGNACIO ROMERO ARAVENA} \\
 	{FECHA: \the\year} \\
 	{PROF. GUÍA: JAVIER RUIZ DEL SOLAR} \\
 	{PROF. CO-GUÍA: PATRICIO LONCOMILLA}
 \end{minipage}
 
 \begin{center}
 	MITIGACIÓN DEL OLVIDO CATASTRÓFICO EN LA CLASIFICACIÓN DE OBJETOS MEDIANTE REPRESENTACIONES GAUSSIANAS DE CLASES
 \end{center}
 
 \noindent Los modelos tradicionales de aprendizaje profundo suelen fallar al trabajar con datos no estacionarios: incorporar clases nuevas exige un reentrenamiento costoso o se olvida lo ya aprendido. Para resolverlo, se presenta un método de aprendizaje incremental por clases libre de gradientes, que aborda el reconocimiento visual continuo mediante la estimación de densidades en forma cerrada, en lugar de la optimización de parámetros. A partir de las representaciones de un extractor de características DINOv2 congelado, cada clase se modela mediante una distribución gaussiana multivariada. Las clases nuevas se integran al instante, sin necesidad de retropropagación, almacenamiento de muestras ni ajuste de parámetros. Como la información de las clases anteriores nunca se altera, el sistema previene el olvido por diseño. Evaluado en Split CIFAR-100, Split ImageNet-R, Split Caltech-256 y en un benchmark conjunto de 30 tareas, este enfoque determinista iguala o supera a los métodos basados en prompts que requieren optimización. La variante MOG4-Shared —que combina una mezcla de gaussianas por clase con una covarianza compartida— ofreció el mejor desempeño. Finalmente, un análisis estadístico del espacio de características de valida el supuesto gaussiano y demuestra que la covarianza ideal depende de la cantidad de datos disponibles por clase.
 %%%%%%%%%%%%%%%%%%%%%%%%%%%%%%%%%%%%%%%
 % SÓLO PARA TESIS EN INGLÉS 
 % BORRAR EN CASO DE ESCRITURA EN ESPAÑOL
 \newpage
 
 \noindent\hfill
 \begin{minipage}[t]{0.6\textwidth}
 	%RESUMEN DE LA TESIS PARA OPTAR AL \\
 	%GRADO DE DOCTOR[A] EN INGENIERÍA \\
 	%ELÉCTRICA, \\
 	RESUMEN DE LA TESIS PARA OPTAR AL \\
 	GRADO DE MAGÍSTER EN CIENCIAS DE LA \\ INGENIERÍA, MENCIÓN ELÉCTRICA, \\
 	MEMORIA PARA OPTAR AL TÍTULO DE \\ INGENIERO CIVIL ELÉCTRICO \\
 	{POR: IGNACIO ROMERO ARAVENA} \\
 	{FECHA: \the\year} \\
 	{PROF. GUÍA: JAVIER RUIZ DEL SOLAR} \\
 	{PROF. CO-GUÍA: PATRICIO LONCOMILLA}
 \end{minipage}
 
 \begin{center}
 	MITIGATING CATASTROPHIC FORGETTING IN  OBJECT CLASSIFICATION VIA GAUSSIAN CLASS REPRESENTATIONS
 \end{center}
 
 \noindent Traditional deep learning models struggle with non-stationary data: adding new categories usually means costly retraining or forgetting what came before. This thesis presents a gradient-free class-incremental learning framework that treats lifelong visual recognition as closed-form density estimation instead of parameter optimization. A frozen DINOv2 encoder is paired with per-class multivariate Gaussian models. New classes are added instantly without backpropagation, sample storage, or tuning. Because past class statistics are never touched, the framework avoids representation drift and forgets nothing by construction. Tested on Split CIFAR-100, Split ImageNet-R, Split Caltech-256, and a 30-task Joint benchmark mixing all three, this deterministic approach matches or beats optimization-heavy, prompt-based methods. MOG4-Shared, combining a small per-class mixture with one shared covariance, is the strongest configuration once cost and consistency are weighed together. A statistical study of the DINOv2 embedding space, using normality tests, spectral analysis, and Bhattacharyya distances, supports the Gaussian assumption behind the method and shows that the best covariance choice depends on how much data each class has.
 

 %En el caso de redactar la tesis en inglés, se debe incluir un resumen adicional en dicho idioma con su título en inglés, ya que el resumen en español (con título en español) es obligatorio independientemente del idioma de la tesis. Ambos resúmenes deben reflejar el mismo contenido. Sin embargo, dado que cada uno debe cumplir independientemente con el límite de doscientas (200) palabras, no es necesario que sean una traducción textual el uno del otro. Las mismas recomendaciones estructurales indicadas en el resumen anterior aplican para ambos.
